\subsection{局部上同调}

在本节中，我们考虑一个局部 Noether 环 A。回忆（VIII, § 3, 第 3 节，引理 2），A 的定义理想就是 A 的异于 A 且包含 \(\mathfrak{m}_{\mathrm{A}}\) 的一个幂的那些理想，或等价地说，是理想 \(\mathfrak{a} \subset \mathfrak{m}_{\mathrm{A}}\) ，它使得 \(\mathrm{A}/\mathfrak{a}\) 为有限长度。我们用 \(\mathcal{D}\) 表示 \(\mathrm{A}\) 的定义理想组成的集合，并赋以与包含关系相反的序关系；它是右有向的。

设 M 为 A-模。使 A 的每个定义理想 a 对应于分次 A-模 \(\operatorname{Ext}_{\mathrm{A}}(\mathrm{A}/\mathfrak{a},\mathrm{M})\) ；若 a 与 b 是满足 \(a \subset b\) 的定义理想，则用 \(p_{ab}: A/a \to A/b\) 表示典范映射，并考虑 A-线性映射 \(\operatorname{Ext}(p_{ab},1_{\mathrm{M}}): \operatorname{Ext}_{\mathrm{A}}(\mathrm{A}/\mathfrak{b},\mathrm{M}) \longrightarrow \operatorname{Ext}_{\mathrm{A}}(\mathrm{A}/\mathfrak{a},\mathrm{M})\) 。这样我们就得到关于有序集 D 的、由分次 A-模与 0 次分次 A-线性映射组成的归纳系统。我们称 M 的局部上同调模，并记之为 \(\mathrm{H}_{\mathrm{A}}(\mathrm{M})\) ，它就是分次 A-模 \(\lim \operatorname{Ext}_{\mathrm{A}}(\mathrm{A}/\mathfrak{a},\mathrm{M})\) 。

理想 \(\mathfrak{m}_{\mathrm{A}}^{n}\)（\(n \geqslant 1\)）构成 \(\mathcal{D}\) 的一个共尾子集；因此有一个分次模的典范同构 \(\varinjlim_{n} \operatorname{Ext}_{\mathrm{A}}(\mathrm{A}/\mathfrak{m}_{\mathrm{A}}^{n}, M) \longrightarrow \mathrm{H}_{\mathrm{A}}(\mathrm{M})\) 。因而，\(\mathrm{H}_{\mathrm{A}}(\mathrm{M})\) 的每个元素都被 \(\mathfrak{m}_{\mathrm{A}}\) 的一个幂零化。

注 1.—— *设 X 为拓扑空间* \(\operatorname{Spec}(A)\) ，\(\mathcal{O}_{\mathrm{X}}\) 为环的结构层，而 \(\widetilde{\mathbf{M}}\) 为与 M 相伴的 \(\mathcal{O}_{\mathrm{X}}\) -模。分次 A-模 \(\mathrm{H}_{\mathrm{A}}(\mathrm{M})\) 等同于模 \(\mathrm{H}_{\{\mathfrak{m}_{\mathrm{A}}\}}(\mathrm{X},\widetilde{\mathrm{M}})\) ，即以闭点 \(\mathfrak{m}_{\mathrm{A}}\) 为支撑的上同调模（该点在 \(X\) 中是闭的）。

对 A-模的每个同态 \(f: \mathrm{M} \to \mathrm{N}\) ，映射 \(\operatorname{Ext}_{\mathrm{A}}(1_{\mathrm{A}/\mathfrak{a}}, f): \operatorname{Ext}_{\mathrm{A}}(\mathrm{A}/\mathfrak{a}, \mathrm{M}) \longrightarrow \operatorname{Ext}_{\mathrm{A}}(\mathrm{A}/\mathfrak{a}, \mathrm{N})\) 构成一个分次线性映射的归纳系统。取归纳极限，得到一个分次同态 \(\mathrm{H}_{\mathrm{A}}(f): \mathrm{H}_{\mathrm{A}}(\mathrm{M}) \to \mathrm{H}_{\mathrm{A}}(\mathrm{N})\) 。对由 A-模与同态组成的每个序列 \(\mathrm{M} \xrightarrow{f} \mathrm{N} \xrightarrow{g} \mathrm{P}\) ，有 \(\mathrm{H}_{\mathrm{A}}(g \circ f) = \mathrm{H}_{\mathrm{A}}(g) \circ \mathrm{H}_{\mathrm{A}}(f)\) 。设

\[
0 \longrightarrow \mathrm{M} \stackrel {f} {\longrightarrow} \mathrm{N} \stackrel {g} {\longrightarrow} \mathrm{P} \longrightarrow 0\tag{ℰ}
\]

为 A-模的一个正合列。由 A, X, 第 90 页，命题 8，扩张模的连接同态 \(\operatorname{Ext}_{\mathrm{A}}(\mathrm{A}/\mathfrak{a},\mathrm{P}) \longrightarrow \operatorname{Ext}_{\mathrm{A}}(\mathrm{A}/\mathfrak{a},\mathrm{M})\) 构成一个 A-线性映射的归纳系统，这些映射是（升）次数 +1 的分次映射。取归纳极限，导出一个 A-同态 \(\partial(\mathcal{E}): \mathrm{H}_{\mathrm{A}}(\mathrm{P}) \to \mathrm{H}_{\mathrm{A}}(\mathrm{M})\) ，它是次数 +1 的分次映射，并使同态序列

\[
\dots \longrightarrow \mathrm{H} _ {\mathrm{A}} ^ {n - 1} (\mathrm{P}) \xrightarrow {\partial^ {n - 1} (\mathcal {E})} \mathrm{H} _ {\mathrm{A}} ^ {n} (\mathrm{M}) \xrightarrow {\mathrm{H} _ {\mathrm{A}} ^ {n} (f)} \mathrm{H} _ {\mathrm{A}} ^ {n} (\mathrm{N}) \xrightarrow {\mathrm{H} _ {\mathrm{A}} ^ {n} (g)} \mathrm{H} _ {\mathrm{A}} ^ {n} (\mathrm{P}) \xrightarrow {\partial^ {n} (\mathcal {E})} \mathrm{H} _ {\mathrm{A}} ^ {n + 1} (\mathrm{M}) \longrightarrow \dots
\]

设 M 为 A-模。对 A 的每个理想 a，A-模 \(\mathrm{Hom}_{\mathrm{A}}(\mathrm{A}/\mathfrak{a},\mathrm{M})\) 典范地等同于 M 中被 a 零化的元素组成的子模。于是 \(\mathrm{H}_{\mathrm{A}}^{0}(\mathrm{M})\) 等同于 M 中被 \(\mathfrak{m}_{\mathrm{A}}\) 的一个幂零化的元素 m 所组成的子模，亦即满足 \(\mathrm{long}_{\mathrm{A}}(\mathrm{Am}) < +\infty\) 的那些元素组成的子模。特别地，当 M 是 Artin 模时，有 \(\mathrm{H}_{\mathrm{A}}^{0}(\mathrm{M}) = \mathrm{M}\) 。

例.—— 1) 若 M 是内射的，则 A-模 \(\mathrm{H}_{\mathrm{A}}^{i}(\mathrm{M})\) 在 \(i > 0\) 时为零，而在 \(i = 0\) 时是内射的（ \(\S 8, \mathrm{n}^{\circ}2\) ，引理 1, c)）。

\begin{enumerate}
\def\labelenumi{\arabic{enumi})}
\setcounter{enumi}{1}
\item
  若 \(\mathrm{H}_{\mathrm{A}}^{0}(\mathrm{M})=\mathrm{M}\) （例如当 M 是 Artin 模时），则 \(\mathrm{H}_{\mathrm{A}}^{i}(\mathrm{M})\) 在 i > 0 时为零。设 (I, e) 为 M 的一个内射包。I 的子模 \(\mathrm{H}_{\mathrm{A}}^{0}(\mathrm{I})\) 是内射的（例 1）且包含 \(e(\mathrm{M})\) ，故等于 I。令 N 为 e 的余核，并考虑正合列 \(0\rightarrow M\xrightarrow{e}\mathrm{I}\xrightarrow{p}\mathrm{N}\rightarrow0\) 。由于 \(\mathrm{I}=\mathrm{H}_{\mathrm{A}}^{0}(\mathrm{I})\) ，有 \(\mathrm{N}=\mathrm{H}_{\mathrm{A}}^{0}(\mathrm{N})\) ，且同态 \(\mathrm{H}_{\mathrm{A}}^{0}(p)\) 是满射。由于 \(\mathrm{H}_{\mathrm{A}}^{i}(\mathrm{I})\) 在 i > 0 时为零（例 1），\(\mathrm{H}_{\mathrm{A}}^{1}(\mathrm{M})\) 为零，且对 i > 1，\(\mathrm{H}_{\mathrm{A}}^{i}(\mathrm{M})\) 同构于 \(\mathrm{H}_{\mathrm{A}}^{i-1}(\mathrm{N})\) ；对整数 i 作归纳即得结论。
\item
  设 \(\Omega\) 为对偶化 A-模。当 \(i \neq \dim(A)\) 时，有 \(\operatorname{Ext}_{A}^{i}(A/a, \Omega) = 0\) ，其中 \(a\) 是 A 的任意一个定义理想（ \(\S 8\) ，\(n^{\circ} 5\) ，例 3），从而 \(\mathrm{H}_{\mathrm{A}}^{i}(\Omega) = 0\) ；当 \(i = \dim(A)\) 时，A-模 \(\mathrm{H}_{\mathrm{A}}^{i}(\Omega)\) 同构于 \(\varinjlim \operatorname{Ext}_{\mathrm{A}}^{i}(\mathrm{A}/\mathfrak{m}_{\mathrm{A}}^{n}, \Omega)\) ，它是 Matlis A-模（同前引，例 6）。
\item
  设 A 为 Noether 局部整环；用 K 表示它的分式域，并设 A ≠ K。K 是内射 A-模（A, X, 第 18 页，例 1），故模 \(\mathrm{H}_{\mathrm{A}}(K)\) 为零（例 1）。由正合列 0 → A → K → K/A → 0，我们对每个 i 导出同构 \(\mathrm{H}_{\mathrm{A}}^{i}(K/A)\) → \(\mathrm{H}_{\mathrm{A}}^{i+1}(A)\)。
\end{enumerate}

更一般地，对每个无扭 A-模 M 及每个整数 i，我们从正合列

\[
0 \rightarrow \mathrm{M} \rightarrow \mathrm{K} \otimes_ {\mathrm{A}} \mathrm{M} \rightarrow (\mathrm{K/A}) \otimes_ {\mathrm{A}} \mathrm{M} \rightarrow 0
\]

导出同构 \(\mathrm{H}_{\mathrm{A}}^{i}((\mathrm{K}/\mathrm{A})\otimes_{\mathrm{A}}\mathrm{M})\to \mathrm{H}_{\mathrm{A}}^{i + 1}(\mathrm{M})\)

\begin{enumerate}
\def\labelenumi{\arabic{enumi})}
\setcounter{enumi}{4}
\tightlist
\item
  保留上例的假设，并再设 dim(A) = 1。设 N 为扭 A-模；由于 A 的每个异于 A 的非零理想都是定义理想（VIII, § 1, 第 3 节，命题 6, e)），有 \(\mathrm{H}_{\mathrm{A}}^{0}(N) = N\)，从而当 i > 0 时 \(\mathrm{H}_{\mathrm{A}}^{i}(N) = 0\) （例 2）。
\end{enumerate}

设 M 为 A-模；用 T(M) 表示它的扭子模。考虑与正合列相伴随的局部上同调长正合列

\[
0 \to \mathrm{T} (\mathrm{M}) \to \mathrm{M} \to \mathrm{M} / \mathrm{T} (\mathrm{M}) \to 0;
\]

考虑到上述结果，由它导出典范同构 T(M) → \(\mathrm{H}_{\mathrm{A}}^{0}(M)\) 与 \(\mathrm{H}_{\mathrm{A}}^{1}(M)\) → \(\mathrm{H}_{\mathrm{A}}^{1}(M/T(M))\)。由于典范同态 \((K/A) \otimes_{A} M\) → \((K/A) \otimes_{A} (M/T(M))\) 是双射，我们最终得到典范同构

\[
\mathrm{H} _ {\mathrm{A}} ^ {0} (\mathrm{M}) \rightarrow \mathrm{T} (\mathrm{M}), \quad \mathrm{H} _ {\mathrm{A}} ^ {1} (\mathrm{M}) \rightarrow (\mathrm{K} / \mathrm{A}) \otimes_ {\mathrm{A}} \mathrm{M}.
\]

命题 1. 设 A 为 Noether 局部环，M 为有限生成的 A-模。

\begin{enumerate}
\def\labelenumi{\alph{enumi})}
\item
  A-模 \(\mathrm{H}_{\mathrm{A}}(\mathrm{M})\) 是 Artin 模，且在次数 \(>\dim (\mathrm{M})\) 处为零。
\item
  令 \(p = \mathrm{prof}_{\mathrm{A}}(\mathrm{M})\) 。有 \(\mathrm{H}_{\mathrm{A}}^{i}(\mathrm{M}) = 0\) （当 \(i < p\) 时），且 \(\mathrm{H}_{\mathrm{A}}^{p}(\mathrm{M}) \neq 0\) （若 \(M\) 非零）。
\end{enumerate}

我们对 \(\dim(M)\) 作归纳来证明 a)。\(\dim(M) \leqslant 0\) 的情形由上文的例 2 得出。设 \(\dim(M) > 0\) ，先取 \(M\) 形如 \(A/p\) ，其中 \(p\) 是 \(A\) 的异于 \(\mathfrak{m}_A\) 的素理想。设 \(x\) 为 \(\mathfrak{m}_A - p\) 的一个元素；有正合列 \(0 \to M \xrightarrow{x_M} M \longrightarrow M/xM \to 0\) ，其中 \(\dim(M/xM) = \dim(M) - 1\) 。由此导出局部上同调的一个正合列

\[
\mathrm{H} _ {\mathrm{A}} ^ {i - 1} (\mathrm{M} / x \mathrm{M}) \longrightarrow \mathrm{H} _ {\mathrm{A}} ^ {i} (\mathrm{M}) \stackrel {x} {\longrightarrow} \mathrm{H} _ {\mathrm{A}} ^ {i} (\mathrm{M}).
\]

\(\mathrm{H}_{\mathrm{A}}^{i}(\mathrm{M})\) 的每个元素都被 \(\mathfrak{m}_{\mathrm{A}}\) 的一个幂零化；因此，要证明该模是 Artin 模，只需证明 \(\mathrm{H}_{\mathrm{A}}^{i}(\mathrm{M})\) 的底座在 \(\kappa_{A}\) 上是有限维的（§ 8, 第 3 节，引理 3）。由归纳假设，\(\mathrm{H}_{\mathrm{A}}^{i}(\mathrm{M})\) 中以 x 为比的乘法映射的核 N 是 Artin 模；由于 x 属于 \(\mathfrak{m}_{\mathrm{A}}\) ，\(\mathrm{H}_{\mathrm{A}}^{i}(\mathrm{M})\) 的底座等同于 N 的底座，因而是有限维的。若 i > dim(M)，则由归纳假设有 \(\mathrm{H}_{\mathrm{A}}^{i-1}(\mathrm{M}/x\mathrm{M}) = 0\) ，故以 x 为比的乘法映射在 \(\mathrm{H}_{\mathrm{A}}^{i}(\mathrm{M})\) 中是单射；由于 \(\mathrm{H}_{\mathrm{A}}^{i}(\mathrm{M})\) 的每个元素都被 x 的一个幂零化，导出 \(\mathrm{H}_{\mathrm{A}}^{i}(\mathrm{M}) = 0\) ，于是在所考虑的情形 a) 成立。

转到一般情形。A-模 M 容许一个合成列 \((\mathrm{M}_{j})_{0\leqslant j\leqslant n}\) ，使每个商 \(M_{j}/M_{j+1}\) 同构于 \(A/p_{j}\) ，其中 \(p_{j}\) 是 A 的素理想（IV, § 1, 第 4 节，定理 1）。对 n 作归纳，证明 M 满足 a)。n=0 的情形是平凡的。正合列 \(0\to M_{1}\to M\to A/p_{0}\to0\) 给出局部上同调的一个正合列

\[
\mathrm{H} _ {\mathrm{A}} ^ {i} \left(\mathrm{M} _ {1}\right) \longrightarrow \mathrm{H} _ {\mathrm{A}} ^ {i} (\mathrm{M}) \longrightarrow \mathrm{H} _ {\mathrm{A}} ^ {i} \left(\mathrm{A} / \mathfrak {p} _ {0}\right);
\]

由归纳假设，A-模 \(\mathrm{H}_{\mathrm{A}}^{i}(\mathrm{M}_{1})\) 是 Artin 模；由已处理的情形，\(\mathrm{H}_{\mathrm{A}}^{i}(\mathrm{A}/\mathfrak{p}_{0})\) 亦然；于是 \(\mathrm{H}_{\mathrm{A}}^{i}(\mathrm{M})\) 是 Artin 模。若 \(i > \dim(\mathrm{M})\) ，模 \(M_{1}\) 与 \(A/p_{0}\) 的维数 < i；从而由归纳假设与已处理的情形，模 \(\mathrm{H}_{\mathrm{A}}^{i}(\mathrm{M}_{1})\) 与 \(\mathrm{H}_{\mathrm{A}}^{i}(\mathrm{A}/\mathfrak{p}_{0})\) 均为零，这蕴含 \(\mathrm{H}_{\mathrm{A}}^{i}(\mathrm{M}) = 0\) 。

设 M 非零，并对整数 \(p = \text{prof}(M)\) 作归纳来证明 b)。\(p = 0\) 的情形由深度的定义得出。设 \(p > 0\) ，选取一个元素 \(x\) （属于 \(\mathfrak{m}_{\mathrm{A}}\) ），使乘法映射 \(x_{\mathrm{M}}\) 是单射。与上面一样，我们得到局部上同调的一个正合列。

\[
\mathrm{H} _ {\mathrm{A}} ^ {i - 1} (\mathrm{M} / x \mathrm{M}) \longrightarrow \mathrm{H} _ {\mathrm{A}} ^ {i} (\mathrm{M}) \xrightarrow {x} \mathrm{H} _ {\mathrm{A}} ^ {i} (\mathrm{M}).
\]

有 \(\operatorname{prof}(M/xM) = \operatorname{prof}(M) - 1\) （ \(\S 1, n^{\circ} 4\) 命题 7），故由归纳假设得 \(H_{A}^{i-1}(M/xM) = 0\) （对 \(i < p\) ），从而与上面一样推出 \(H_{A}^{i}(M) = 0\) 。特别地，\(H_{A}^{p-1}(M)\) 为零，于是同态 \(H_{A}^{p-1}(M/xM) \longrightarrow H_{A}^{p}(M)\) 是单射；从而由归纳假设，\(H_{A}^{p}(M)\) 非零。

可以证明，当 M 非零时，模 \(H_{A}^{\dim(M)}(M)\) 非零（习题 4；参见第 3 节定理 2 的推论）。

推论. 设 M 为非零的有限生成 Macaulay A-模。A-模 \(\mathrm{H}_{\mathrm{A}}^{i}(\mathrm{M})\) 在 \(i \neq \dim(\mathrm{M})\) 时为零，而在 \(i = \dim(\mathrm{M})\) 时非零。

注 2.—— 对 A 的每个定义理想 a，A-模 \(\operatorname{Ext}_{\mathbf{A}}(\mathbf{A}/\mathfrak{a},\mathbf{M})\) 被 a 零化，且 A/a 等同于 \(\widehat{\mathbf{A}}/\mathfrak{a}\widehat{\mathbf{A}}\) ；因此，分次 A-模 \(\operatorname{Ext}_{\mathbf{A}}(\mathbf{A}/\mathfrak{a},\mathbf{M})\) 等同于 \(\widehat{\mathbf{A}}\otimes_{\mathbf{A}}\operatorname{Ext}_{\mathbf{A}}(\mathbf{A}/\mathfrak{a},\mathbf{M})\) ，从而也等同于 \(\operatorname{Ext}_{\widehat{\mathbf{A}}}(\widehat{\mathbf{A}}/\mathfrak{a}\widehat{\mathbf{A}},\widehat{\mathrm{A}}\otimes_{\mathbf{A}}\mathbf{M})\) （A, X, 第 111 页，命题 10）。诸理想 \(\mathfrak{a}\widehat{\mathbf{A}}\) （其中 \(\mathfrak{a}\in\mathcal{D}\) ）组成的集合包含 \(\mathfrak{m}_{\widehat{\mathbf{A}}}\) 的诸幂，故在 \(\widehat{\mathbf{A}}\) 的定义理想组成的集合中共尾；于是由前述得到分次 A-模的一个同构

\[
\mathrm{H} _ {\mathrm{A}} (\mathrm{M}) \longrightarrow \mathrm{H} _ {\widehat {\mathrm{A}}} (\widehat {\mathrm{A}} \otimes_ {\mathrm{A}} \mathrm{M}).
\]

若 A-模 M 是有限型的，则 A-模 \(\widehat{\mathrm{A}}\otimes_{\mathrm{A}}\mathrm{M}\) 等同于 M 的完备化 \(\widehat{\mathrm{M}}\) （III, § 3, 第 4 节，定理 3），且有一个次数为 0 的分次同构 \(\mathrm{H}_{\mathrm{A}}(\mathrm{M})\to \mathrm{H}_{\widehat{\mathrm{A}}}(\widehat{\mathrm{M}})\) 。

